\chapter{电压与电压表}\label{ch:g8:voltage-voltmeter}

自第一块电池点亮第一只灯泡起，本书中便有一个词承担着
沉重而含糊的任务：“推”。电池推得轻，市电推得凶，
两块首尾相接的电池合力而推。
“推”如今成为一个可测量的量，
拥有著名的单位，而它的仪器，与上一章的不同，
拒绝站在路中间。

\section{电压}

\begin{definition}[电压]\label{def:g8:voltage-voltmeter:voltage}
电路中两点之间的\emph{电压}\index{电压} $U$，
度量两点之间可用的电气“推力”——
即电流从一点被推向另一点的力度，
正如瀑布的落差度量它推动水轮的力度。
电压总是\emph{两点之间}的量：电池的两个接线柱之间，灯泡的两端之间，
绝不只限于某一个点。
其单位是
\emph{伏特}\index{伏特}（\unit{V}），以纪念第一块电池的
发明者。
\end{definition}

\begin{example}[伏特家族相册]\label{ex:g8:voltage-voltmeter:values}
单节圆形电池：\qty{1.5}{V}。扁平的袖珍电池：
\qty{4.5}{V}——外包装内三节电池首尾相接。
长方形九伏电池，两个接线柱都在顶部：
\qty{9}{V}，内含六节隐藏的电芯。汽车电池：\qty{12}{V}。墙上的
插座：\qty{230}{V}——比一百五十多节小电池还要多的“推力”，
这正是“插座，绝不可碰”这句话的全部定量内容。
云与地之间的闪电：
高达数亿伏特，但转瞬即逝。
\end{example}

\begin{proposition}[无电流时的电压]\label{prop:g8:voltage-voltmeter:nocurrent}
电压与电流是两个不同的量，而且电压是先行的一方。
一节单独放在抽屉里、不接任何导线的电池，
不会驱动任何电流——但它两端之间
完整的\qty{1.5}{V}电压依然存在，可以测量，
就像水磨停工的那天，瀑布的落差依旧。
推力可以存在而无水流；水流却永远离不开推力。
\end{proposition}

\begin{definition}[电压表]\label{def:g8:voltage-voltmeter:voltmeter}
\emph{电压表}\index{电压表}测量两点之间的电压——
因此它应当\emph{跨接}在这两点之间：
与所研究的元件\emph{并联}，两根表笔分别接在两侧，
好比测量员把水准仪架在两个高度之间。
与电流表在各方面都截然相反，它几乎
\emph{完全}地阻碍电流：跨接在任何元件上，它几乎分不走
一点电流，因此接上它不会扰动任何东西，
也无需断开任何回路。
\end{definition}

\begin{omfigure}
\begin{tikzpicture}[scale=0.9]
  \draw (0,0) to[battery1, l=battery] (0,2.8)
    to[short] (2.0,2.8)
    to[lamp, l=lamp] (4.6,2.8)
    to[short] (6.4,2.8)
    to[short] (6.4,0)
    to[short] (0,0);
  \draw (2.0,2.8) to[short] (2.0,4.6)
    to[voltmeter, l=$U$] (4.6,4.6)
    to[short] (4.6,2.8);
\end{tikzpicture}

{\small 电压表采取测量员的姿态：跨接在
灯泡两侧，每侧一根表笔，测量两端之间的“推力”——
绝不站在路中间。}
\end{omfigure}

\begin{method}[测量电压]\label{met:g8:voltage-voltmeter:measure}
\begin{enumerate}
  \item 保持电路完好——对电压表而言，
        无须断开任何部分；
  \item 将一根表笔接在元件的每一侧：
        \unit{V}表笔接在常规电流流入的一端，
        公共端接在电流流出的一端；
  \item 先用最大量程，再逐档调至合适量程；
  \item 读数，并写出单位。
\end{enumerate}
两表的“脾气”并列如下：电流表——断开电路，
站入其中，不设阻碍；
电压表——不断开电路，站在一旁，处处阻碍电流。混淆两者的
“脾气”，是实验室里代价最高的习惯。
\end{method}

\section{额定电压}

\begin{definition}[额定电压]\label{def:g8:voltage-voltmeter:nominal}
每个灯泡、电动机和电器都标有\emph{额定
电压}\index{额定电压}——即设备设计时所针对的电压，
印在底座或铭牌上：手电筒灯泡“3.5 V”、汽车前灯“12 V”、
水壶“230 V”。以额定电压供电时，
设备按设计运行；供电远低于此，它便闹脾气——
灯泡昏暗，电动机无力；供电远高于此，它便烧毁——
亮得惊人，却只此一次。
\end{definition}

\begin{example}[灯泡与电池的匹配]\label{ex:g8:voltage-voltmeter:matching}
标有“3.5 V”的灯泡接在\qty{4.5}{V}扁平电池上：略微过载，
灯光明亮，但寿命短暂。接在单节
\qty{1.5}{V}电池上：只有昏昏欲睡的橙色微光。接在\qty{9}{V}
电池上：刹那间耀眼，随后细丝断裂——
这是灯丝的讣告。标签即契约：遵守它，
灯泡便能活出标称的寿命。
\end{example}

\begin{remark}[为何 230 V 会致命]\label{rem:g8:voltage-voltmeter:mains}
现在，那条古老的绝对规则有了它的数字。人体潮湿的
导电通路，接上\qty{4.5}{V}时，通过的电流
小得无法察觉。同一通路接上\qty{230}{V}——电压是原来的五十倍——
通过的电流也是五十倍，远超足以使心脏停跳的
百分之几安培的电流。插座并非
“只是略强”于你的电池实验台：它比之高出
两个数量级，而规章的严厉程度，
不过是简单的算术而已。
\end{remark}

\section{初步测量电压}

\begin{example}[绕简单回路一周]\label{ex:g8:voltage-voltmeter:readings}
一节电池、一只灯泡，手持电表。测量电池两端：
\qty{4.5}{V}。测量灯泡两端：\qty{4.5}{V}——灯泡
得到的正是电池提供的电压。测量一段
连接\emph{导线}两端：\qty{0}{V}，或电表所能示出的最接近的值——
示数。好导线如同平坦的道路，不耗费任何“推力”。
“推力”存在于\emph{用电器}两端，而不在导线上：
这一读数值得记住，待下一章的定律将它
正式确认。
\end{example}

\begin{example}[断开开关的“供认”]\label{ex:g8:voltage-voltmeter:openswitch}
灯泡不亮，开关断开——推力去了哪里？测量：
不亮的灯泡两端为 \qty{0}{V}；两端的\emph{断开的
开关}则为满额 \qty{4.5}{V}。电池的推力完整地
横跨缺口等待——回路中唯一一处
两点在电学上真正“相距甚远”。合上
开关，两处读数互换。断开的开关
正是电压等待的地方。
\end{example}

\section{练习}

\begin{exercise}[$\star$]\label{exo:g8:voltage-voltmeter:1}
电压测量的是什么？为什么它总是“两点之间”的量？
写出电压的单位，并说明水流比喻。
\end{exercise}

\begin{exercise}[$\star$]\label{exo:g8:voltage-voltmeter:2}
凭记忆写出：圆柱形电池、扁平电池、
方形电池、汽车电池和市电的电压。
\end{exercise}

\begin{exercise}[$\star$]\label{exo:g8:voltage-voltmeter:3}
一节电池未接任何用电器。有多大电流通过？
跨接其两端的电压表示数是多少？哪一条命题在此适用？
\end{exercise}

\begin{exercise}[$\star$]\label{exo:g8:voltage-voltmeter:4}
比较两种电表的用法：连接方式、是否“断路”、
对电流的阻碍作用。哪一种错误会借助表盘造成
短路？
\end{exercise}

\begin{exercise}[$\star$]\label{exo:g8:voltage-voltmeter:5}
什么是额定电压？预测一个“6 V”灯泡
在 \qty{1.5}{V}、\qty{6}{V}、\qty{12}{V} 下的表现。
\end{exercise}

\begin{exercise}[$\star$]\label{exo:g8:voltage-voltmeter:6}
在简单回路中，电压表跨接电池、灯泡
以及一段导线两端时，分别读数多少？推力
存在于何处？
\end{exercise}

\begin{exercise}[$\star$]\label{exo:g8:voltage-voltmeter:7}
开关断开时，不亮的灯泡两端和断开的开关两端各是多少？
合上开关的瞬间，两处读数会怎样变化？
\end{exercise}

\begin{exercise}[$\star$]\label{exo:g8:voltage-voltmeter:8}
为何把电压表并接在正在工作的灯泡两端，几乎不会扰动
电路？这只电压表拒绝
去做什么？
\end{exercise}

\begin{exercise}[$\star\star$]\label{exo:g8:voltage-voltmeter:9}
扁平的 \qty{4.5}{V} 电池内藏三节单体电池。画出（或描述）
它们的排列方式，并计算若其中一节被误装反时的电压
——即多年前手电筒一章的内容，如今用数字来算。
（方向相反的推力相减。）
\end{exercise}

\begin{exercise}[$\star\star$]\label{exo:g8:voltage-voltmeter:10}
露营者的灯泡额定 \qty{12}{V}，房车电池新充满时读数为
\qty{12.6}{V}，经过一个漫长的夜晚后为 \qty{11.2}{V}。
描述灯泡的这一晚，并说明它为何没有烧坏。
\end{exercise}

\begin{exercise}[$\star\star$]\label{exo:g8:voltage-voltmeter:11}
误将电压表与灯泡\emph{串联}接入（回路断开，电表接在缺口处）：
预测灯泡的表现
与电表的示数——别忘了电压表所阻碍的是什么。
（灯泡不亮与电表示数都很有启发意义，
值得细想。）
\end{exercise}

\begin{exercise}[$\star\star\star$]\label{exo:g8:voltage-voltmeter:12}
两只相同的灯泡串联接在 \qty{4.5}{V} 的电源上：电压表
测得每只灯泡两端约为 \qty{2.25}{V}。将这两只灯泡
并联接在同一电池上：预测每只灯泡的电压，并
用你多年来已知的亮度规律调和两种结果
（串联的一对较暗，并联的一对
全亮）。哪一章将使你的预测成为定律？
\end{exercise}

\section{问题：电池抽屉大检查}

\begin{problem}[{周末问题——电池抽屉的
大检查；一只电压表、一盒散装单体电池，以及家中的
小电器架}]\label{pb:g8:voltage-voltmeter:1}
家家都有这样一个抽屉：散落的单体电池、来历不明的电池，
以及“需要换新”的小电器。检查日——电压表
在手。

\medskip\noindent\textbf{Part I --- 单体电池分类。}
\begin{enumerate}
  \item 描述用电压表检查一节
        散放的 \qty{1.5}{V} 单体电池时的正确连接方式
        ——并说明为何无需先搭建电路。
  \item 抽屉中刚取出的单体电池读数如下：
        \qty{1.58}{V}、\qty{1.32}{V}、\qty{0.9}{V}、
        \qty{0.1}{V}。请将其分类：新鲜、可用、电量偏低、已耗尽。
  \item 一节九伏电池的读数为 \qty{8.9}{V}。大致
        每节内部单体电池贡献多少？
  \item 为何一节将耗尽的单体电池仍显示\emph{一定}
        电压，尽管它已无法点亮任何东西？（它的
        推力仍微弱地存在；它失去了何种能力，无法
        驱动？）
\end{enumerate}

\medskip\noindent\textbf{第二部分 --- 小物件架。}
\begin{enumerate}[resume]
  \item 相机需要两节电池首尾相接，它需要多大电压，
        以及第 2 题中剩下的那两节
        \qty{1.32}{V} 电池能提供多少？
  \item 旧收音机要用四节电池首尾相接，
        电压低于 \qty{5}{V} 就罢工。抽屉里的电池中，
        哪些（如果有的话）能让它工作？
  \item 弟弟把三节电池硬塞进只能装两节的
        手电筒里，说是“为了更亮”，还把第三节横卡在
        弹簧上。灯泡额定电压为 \qty{3}{V}。
        请用额定电压的语言预测会出现什么情况。
  \item 门铃变压器的铭牌上写着“\qty{8}{V}”。
        有人提议把门铃直接接到一节九伏电池上，
        理由是“八和九差不多”。
        请评估：这是合理搭配，还是在慢慢损坏门铃？
\end{enumerate}

\medskip\noindent\textbf{第三部分 --- 审计定律。}
\begin{enumerate}[resume]
  \item 在手电筒中（灯泡与电池串联，开关断开），
        审计员测得：断开的开关两端
        电压为 \qty{3.1}{V}；灯泡两端电压为 \qty{0}{V}。请解释两个读数，
        并预测开关闭合（咔嗒一声）后两者各为多少。
  \item 在手电筒线路的每一段两端，
        电表读数均为 \qty{0}{V}，无论灯亮还是灯灭。
        用“推力”的语言说明，这对良好导线意味着什么？
  \item 按家庭的既定规矩，市电插座\emph{不}在审计之列。
        请用本章给出的数值重新表述这一定律，
        并写出其背后的算式。
  \item 写审计报告，共三句 --- 电压表是什么、又绝不是什么；
        标签上额定电压承诺了什么、又威胁了什么；
        以及在任何闲置电路中，电压在何处等待。
        ——
\end{enumerate}
\end{problem}
